% SPDX-FileCopyrightText: 2026 Will Cook
% SPDX-License-Identifier: CC-BY-4.0
%
% One of the Erdős Problem Notes: a short standalone treatment of a single
% problem in the ErdosProblems expansion library.  Build with
% `tectonic erdos-243-reciprocal-tail-rigidity.tex` or `make`.
\documentclass[11pt]{article}
% SPDX-FileCopyrightText: 2026 Will Cook
% SPDX-License-Identifier: CC-BY-4.0
%
% Shared preamble for the Erdős Problem Notes: the short, standalone,
% problem-owned manuscripts covering the expansion library `ErdosProblems`.
%
% One note owns one Erdős problem.  Each note repeats whatever context its own
% reader needs, so that a reader who arrives at a single problem never has to
% assemble it from the gateway paper.  Deliberate repetition across notes is the
% design, not an oversight.
%
% Source links resolve against one pinned formal-source commit, declared once
% below.  `scripts/check_problem_note_sources.py` verifies that every authored
% (file, line, declaration) triple names that exact declaration in the pinned
% snapshot, so a link cannot silently rot as later work moves lines.

\usepackage{paper-house-style}
\usepackage{array}
\usepackage{booktabs}
\usepackage{enumitem}
\setlist{topsep=0.35em,itemsep=0.2em}
\usepackage[colorlinks=true,linkcolor=linkinternal,urlcolor=linkexternal,
  citecolor=linkinternal]{hyperref}
\hypersetup{bookmarksnumbered=true,bookmarksopen=true,bookmarksopenlevel=1,
  pdfauthor={Will Cook},pdflang={en-GB}}

% ---------------------------------------------------------------------------
% Pinned formal source.  \commit names the checkpoint every link resolves
% against; the checker reads that snapshot rather than the working tree, so the
% printed coordinates stay correct however later waves move declarations.
% ---------------------------------------------------------------------------
\newcommand{\commit}{e5fd26a6d1b1a346430205eab5385b4c9565d004}
\newcommand{\commitshort}{e5fd26a6d1b1}
\newcommand{\repobase}{https://github.com/wcook04/plectis-lean-erdos249-257/blob/\commit}
\newcommand{\sourceurl}{https://github.com/wcook04/plectis-lean-erdos249-257/tree/\commit}
\newcommand{\PX}{ErdosProblems}

% Declaration names stay inside link targets.  The printed label is either a
% spaced, lower-case reading of the declaration name (\lref) or a short authored
% phrase (\lword); neither prints a module path or a raw Lean identifier.
\ExplSyntaxOn
\cs_new_protected:Npn \noteleanlabel #1
  {
    \tl_set:Nx \l_tmpa_tl { \tl_to_str:n {#1} }
    \regex_replace_all:nnN { _+ } { \c{space} } \l_tmpa_tl
    \regex_replace_all:nnN
      { ([a-z0-9])([A-Z]) }
      { \1\c{space}\2 }
      \l_tmpa_tl
    \tl_set:Nx \l_tmpa_tl { \tl_lower_case:n { \tl_use:N \l_tmpa_tl } }
    \tl_use:N \l_tmpa_tl
  }
\ExplSyntaxOff
\newcommand{\leanlabel}[1]{\noteleanlabel{#1}}
\newcommand{\lref}[3]{\href{\repobase/\PX/#1\#L#2}{\leanlabel{#3}}}
\newcommand{\lword}[4]{\href{\repobase/\PX/#1\#L#2}{#4}}
\newcommand{\lloc}[2]{\href{\repobase/\PX/#1\#L#2}{Lean source}}

% Links into a sibling library, named repository-relative.  The expansion
% library is implicit in \lref/\lword, because that is where problem-owned
% work lands.  Several problems have their headline theorem in the older
% reviewed #249/#257 corpus, and a note that could not name that library
% would have to paraphrase a checked statement instead of linking it.
\newcommand{\mref}[3]{\href{\repobase/#1\#L#2}{\leanlabel{#3}}}
\newcommand{\mword}[4]{\href{\repobase/#1\#L#2}{#4}}
\newcommand{\mloc}[2]{\href{\repobase/#1\#L#2}{Lean source}}
\emergencystretch=2em

\newtheorem{theorem}{Theorem}[section]
\newtheorem{proposition}[theorem]{Proposition}
\newtheorem{lemma}[theorem]{Lemma}
\newtheorem{corollary}[theorem]{Corollary}
\theoremstyle{definition}
\newtheorem{definition}[theorem]{Definition}
\newtheorem{example}[theorem]{Example}
\newtheorem{problem}[theorem]{Problem}
\theoremstyle{remark}
\newtheorem*{remark}{Remark}

\newcommand{\N}{\mathbb{N}}
\newcommand{\Npos}{\mathbb{N}_{>0}}
\newcommand{\Z}{\mathbb{Z}}
\newcommand{\Q}{\mathbb{Q}}
\newcommand{\R}{\mathbb{R}}
\newcommand{\lcm}{\operatorname{lcm}}
\newcommand{\ph}{\varphi}

% The series line printed under every note's title block.
\newcommand{\seriesline}{Erd\H{o}s Problem Notes}

% Production provenance.  The wording is fixed by the corpus provenance
% contract and reproduced verbatim in every manuscript in this repository, so
% the papers cannot drift into different accounts of how they were made.
\newcommand{\noteprovenance}{\provenancenote{\emph{Authorship and AI use.} The
prose, formal proofs, and software in this version were drafted and revised by
large-language-model agents under Will Cook's direction.  Cook set the
objectives and acceptance criteria, reviewed and authorised the manuscript, and
is responsible for its contents.  The agents are production tools, not authors.
See \emph{Statements and declarations}.}}

% The status paragraph every note carries, in its own words, before any result.
% Kernel checking establishes the exact Lean proposition; it does not establish
% that the proposition is the interesting one, that it is new, or that the
% Erdős problem is closed.
\newcommand{\statusboundary}{%
  \emph{Status.}  The problem treated here is open, and this note does not
  close it.  Every statement below marked as checked is a proposition that the
  pinned Lean kernel accepts from the sources this note links to, with no
  \texttt{sorry}, no added axiom, and no unchecked evaluation.  That is a claim
  about the formal statement, not about its mathematical interest, its
  novelty, or the original problem.  The unresolved obligations are named
  exactly, in their own section, and none of the finite computations,
  reductions, or no-go results here removes one of them.}

% Problem-specific source pin: #243's cancellation and feedback interfaces
% advance independently of the corpus default.
\renewcommand{\commit}{32545d77c4b5cfd72ebd36c8dcd2418bf7cc4d31}
\renewcommand{\commitshort}{32545d77c4b5}
\hypersetup{
  pdftitle={Excluding the bounded negative part: Lean-checked rigidity for Erdős Problem 243},
  pdfsubject={Erdős Problem Notes: Problem 243},
  pdfkeywords={irrationality, Ahmes series, Sylvester sequence, unit fractions, Lean 4},
}

% Names for the state maps and the two search objects.  They are set as
% upright operators so that they read as maps rather than as products.
\newcommand{\sylnext}{\operatorname{syl}}
\newcommand{\den}{\operatorname{den}}
\newcommand{\tail}{\operatorname{tail}}
\newcommand{\ctr}{\operatorname{ctr}}
\newcommand{\defect}{\operatorname{def}}
\newcommand{\fnum}{\operatorname{num}}
\newcommand{\surv}{\operatorname{Surv}}

\runninghead{Erd\H{o}s \#243: excluding the bounded negative part}
\title{Excluding the Bounded Negative Part}
\subtitle{Lean-checked rigidity for Erd\H{o}s Problem~\#243}
\author{Will Cook}
\date{July 2026}

\begin{document}
\maketitle
\noteprovenance

\begin{abstract}
Let $a_1<a_2<\cdots$ be positive integers with $a_{n+1}/a_n^2\to1$
and rational reciprocal sum, and put $P_n=\prod_{j<n}a_j$. We prove that
\[
 \limsup_{n\to\infty}\frac{P_n}{a_n}
 \left(\frac{a_n^2}{a_{n+1}}-1\right)<+\infty
\]
forces $a_{n+1}=a_n^2-a_n+1$ eventually. In the integral tail state,
bounded negative error caps upward increments, while normalised vanishing
forces a nonterminal numerator to infinity. Stabilising the common divisor
then supplies a Chinese remainder block that those increments cannot cross.
An LCM formulation gives a weighted record criterion and the corresponding
LCM-prefactor consequence. The required bound or record budget remains
unproved for the unrestricted problem.
\end{abstract}

\section{Introduction}\label{sec:problem}

\begin{corollary}[original-coordinate bounded defect]\label{res:originalbounded}
Let $a_1<a_2<\cdots$ be positive integers, $a_{n+1}/a_n^2\to1$, and
$\sum_{n\ge1}1/a_n\in\Q$.  Put $P_n=\prod_{j<n}a_j$.  If
\[
 \limsup_{n\to\infty}\frac{P_n}{a_n}
 \left(\frac{a_n^2}{a_{n+1}}-1\right)<+\infty,
\]
then $a_{n+1}=a_n^2-a_n+1$ for all sufficiently large $n$.
\end{corollary}

\paragraph{The obstruction.}
The equation $C_{n+1}=C_n-E_n$ translates a bound on negative error into a
bound on upward increments. At negative indices the common divisor of
$C_n,D_n$ divides the bounded magnitude. Once this divisor stabilises,
reduction makes the later numerators coprime to the old multipliers.
The Chinese remainder theorem places consecutive forbidden heights beyond
any fixed prefix. A first crossing then gives the contradiction.
Section~\ref{sec:lcmrecords} uses the same divisibility at record sources,
so that only record excess is charged.

\begin{samepage}
Write $\ctr(a,D,C)=D-(a-1)C$. The following integer-state theorem
implies Corollary~\ref{res:originalbounded}; the ordinary tail transfer is
proved in Section~\ref{sec:transfer}. Its centring hypothesis is redundant
under normalised vanishing and is retained to match the source statement.

\begin{theorem}[bounded negative part]\label{res:bounded}
Let $a,C,D:\N\to\N$ and $E:\N\to\Z$ satisfy
\begin{enumerate}
\item $a_n>1$ and $C_n>0$ for every $n$;
\item the exact dynamics $C_{n+1}+D_n=a_nC_n$ and $D_{n+1}=a_nD_n$;
\item $E_n=\ctr(a_n,D_n,C_n)$ for every $n$;
\item \emph{eventual strict centring}: $|E_n|<C_n$ for all large $n$;
\item \emph{eventually bounded negative part}: $-B\le E_n$ for all large $n$,
for some $B$;
\item \emph{normalised vanishing}: for every $K$ there is an $N$ with
$K\,|E_n|<C_n$ for all $n\ge N$.
\end{enumerate}
Then $E_n=0$ for all sufficiently large $n$.
\end{theorem}
\end{samepage}

\begin{problem}[Erd\H{o}s \#243]\label{res:problem}
Let $1\le a_1<a_2<\cdots$ be a sequence of integers with
\[
 \lim_{n\to\infty}\frac{a_n}{a_{n-1}^{2}}=1
 \qquad\text{and}\qquad
 \sum\frac{1}{a_n}\in\Q .
\]
Then $a_n=a_{n-1}^{2}-a_{n-1}+1$ for all sufficiently large $n$.
\end{problem}

The additional hypothesis in Theorem~\ref{res:bounded} is the lower bound
on $E_n$. The growth and rationality assumptions already supply the other
state hypotheses after a finite shift. Koizumi's nonnegative-error case is
integer descent \cite[Proposition~1(2)]{koizumi2025}; the theorem allows
arbitrary sign changes with bounded negative depth. Bado's author-posted
preprint assumes two-sided bounded errors \cite[Theorem~5.1]{bado2026}.
The comparisons concern these stated hypotheses and make no priority claim.

\section{Proof of bounded-negative rigidity}\label{sec:bounded}
The arithmetic ingredients are proved immediately afterwards. Normalised
vanishing and integrality give $C_n\to\infty$ on a nonzero tail, while the
lower error bound gives $C_{n+1}\le C_n+B$.

\begin{proof}[Proof of Theorem~\ref{res:bounded}]
Suppose not.  We first remove every late zero: by Theorem~\ref{res:absorb} and
its contrapositive form, $E$ is nowhere zero beyond the centring threshold.
Two cases remain.  If $E$ is negative cofinally,
then~(5) bounds those magnitudes, so Proposition~\ref{res:gcdstab} makes the
tail gcd stabilise and the orbit may be taken reduced past that point.
Hypothesis~(5) with
Proposition~\ref{res:update} gives $C_{n+1}\le C_n+B$, a bounded rise;
hypothesis~(6) with $E$ nowhere zero gives $C_n\to\infty$.  A reduced exact
tail with a divergent numerator and bounded rise is excluded by
Theorem~\ref{res:barrier}.  It remains to treat the case that $E$ is eventually
positive, and there the descent
of Theorem~\ref{res:descent} applies and forces $E$ to vanish, contradicting
nowhere-vanishing.
\end{proof}

\begin{corollary}\label{res:cor}
Under Theorem~\ref{res:bounded}, the multipliers satisfy
$a_{n+1}=a_n^2-a_n+1$ eventually.
\end{corollary}
\begin{proof}
Apply the two-zero implication in Corollary~\ref{res:eventual} below.
\end{proof}

\section{The arithmetic ingredients}\label{sec:state}
\subsection{Error, absorption and descent}\label{sec:defect}\label{sec:descent}
Write $\Delta_n=a_{n+1}-(a_n^2-a_n+1)$.

\begin{proposition}[error identities]\label{res:update}\label{res:defect}
For an exact integer state,
\[
 C_{n+1}=C_n-E_n,\qquad
 \Delta_n C_{n+1}=a_n^2E_n-E_{n+1}.
\]
\end{proposition}
\begin{proof}
Substitute $E_n=D_n-(a_n-1)C_n$ and use
$D_{n+1}=a_nD_n$, $C_{n+1}=a_nC_n-D_n$.
\end{proof}

\begin{theorem}[absorption and descent]\label{res:absorb}\label{res:descent}
For a positive exact state with strict centring, $E_n=0$ implies
$E_{n+1}=0$. For any positive integer state with $C_{n+1}=C_n-E_n$,
eventual nonnegativity of $E_n$ implies its eventual vanishing.
\end{theorem}
\begin{proof}
When $E_n=0$, the second identity makes $E_{n+1}$ a multiple of
$C_{n+1}$; strict centring forces that multiple to be zero. In the second
assertion, $C_n$ is eventually a nonincreasing sequence of positive
integers, so it stabilises.
\end{proof}

\begin{corollary}[two zero errors]\label{res:step}\label{res:eventual}
If $E_n=E_{n+1}=0$ and $C_{n+1}\ne0$, then
$a_{n+1}=a_n^2-a_n+1$. Thus eventual zero error in a positive exact state
implies the eventual Sylvester recurrence.
\end{corollary}
\begin{proof}
The second error identity has nonzero factor $C_{n+1}$.
\end{proof}

Absorption holds after any eventual centring threshold. On a nonterminal
tail it removes all subsequent zeros. Then $|E_n|\ge1$, so normalised
vanishing gives $C_n\to\infty$.

\subsection{The first-crossing obstruction}\label{sec:barrier}
\begin{lemma}[consecutive multiples]\label{res:crt}
For pairwise coprime integers $m_0,\ldots,m_{B-1}\ge2$ and every lower
bound, there is a larger $t$ such that $m_i\mid t+i$ for each $i<B$.
\end{lemma}
\begin{proof}
Solve $t\equiv-i\pmod{m_i}$ by the Chinese remainder theorem, then add
multiples of $\prod_i m_i$.
\end{proof}

\begin{theorem}[bounded-rise obstruction]\label{res:barrier}
Let $u_n\in\N$ tend to infinity with $u_{n+1}\le u_n+B$, for a fixed
integer $B>0$. There is no infinite pairwise coprime family $m_i\ge2$
with $\gcd(m_i,u_t)=1$ whenever $i<t$.
\end{theorem}
\begin{proof}
Choose $B$ old moduli and take a CRT block $[t,t+B)$ beyond every numerator
in the prefix before they become old. At the first crossing of $t$, the
new numerator lies in that block, so it is divisible by an old modulus.
This contradicts coprimality. No monotonicity of $u$ is used.
\end{proof}

\subsection{Reduction and stabilisation}
A reduced exact tail has positive $u_n$, integer $v_n\ge0$,
$\gcd(u_n,v_n)=1$, and
\[
 u_{n+1}=a_nu_n-v_n,\qquad v_{n+1}=a_nv_n.
\]
\begin{proposition}[persistent coprimality]\label{res:reduced}
In a reduced exact tail, $\gcd(a_n,v_n)=1$. Distinct multipliers are
pairwise coprime, and every earlier multiplier is coprime to every later
numerator.
\end{proposition}
\begin{proof}
A common prime divisor of $a_n,v_n$ would divide both $u_{n+1}$ and
$v_{n+1}$. Also, $a_i\mid v_t$ for $i<t$, so reducedness gives
$\gcd(a_i,u_t)=1$, and $\gcd(a_t,v_t)=1$ gives
$\gcd(a_i,a_t)=1$.
\end{proof}

\begin{proposition}[gcd stabilisation]\label{res:gcdstab}
For a positive exact state, if negative errors have bounded magnitudes
along a cofinal set of indices, then $G_n=\gcd(C_n,D_n)$ eventually
stabilises. Division by its stable value gives a reduced exact tail.
\end{proposition}
\begin{proof}
Both updates preserve common divisors, so $G_n\mid G_{n+1}$. Moreover
$G_n\mid E_n$, hence $G_n\le -E_n$ at a negative index. The positive
divisibility chain is bounded along a cofinal set and therefore bounded
everywhere; it eventually stabilises. Division by the stable value
preserves both exact updates and leaves the states coprime.
\end{proof}

In the main proof, division by the stable gcd preserves divergence and a
bounded upward increment. Theorem~\ref{res:barrier} therefore applies to
the reduced numerator. Sparse changes of the gcd without eventual
stabilisation are a different statement; the supporting record retains
that distinction and its finite-block consequences.

\section{Transfer to a reciprocal series}\label{sec:transfer}
Let $x_n=\sum_{k\ge n}1/a_k$, suppose $x_1=p/q$ with $q>0$, and put
\[
 P_n=\prod_{j<n}a_j,\qquad D_n=qP_n,\qquad C_n=D_nx_n.
\]
Each $C_n$ is a positive integer, since $D_n$ clears the rational sum and
every term of the preceding finite prefix. The tail relation
$x_n=1/a_n+x_{n+1}$ gives the exact state updates.

Quadratic growth gives $a_n\ge\exp(c2^n)$ eventually for some $c>0$ and
\[
 x_n=\frac1{a_n}+\frac1{a_{n+1}}+O(a_{n+1}^{-2}).
\]
Consequently
\[
 \frac{C_{n+1}}{C_n}
 =\frac{a_nx_{n+1}}{x_n}
 =\frac{a_n^2}{a_{n+1}}+O(1/a_n)\longrightarrow1,
 \qquad \frac{E_n}{C_n}\longrightarrow0.
\]
Thus strict centring also holds eventually. These are the ordinary
canonical-tail implications of Koizumi's Corollary~3 and Lemma~4
\cite{koizumi2025}.

\begin{proof}[Proof of Corollary~\ref{res:originalbounded}]
Put $\gamma_n=a_n^2/a_{n+1}-1$. Since
\[
 \frac{P_{n+1}/a_{n+1}}{P_n/a_n}=1+\gamma_n\longrightarrow1,
 \qquad P_n/a_n=\exp(o(n)),
\]
the two-term estimate gives
\begin{equation}\label{eq:canonical-dictionary}
 E_n+q\frac{P_n}{a_n}\gamma_n
 =\frac{qP_n}{a_{n+1}}+
 O\!\left(\frac{qP_na_n}{a_{n+1}^2}\right)=o(1).
\end{equation}
The assumed upper bound therefore gives an eventual lower bound on the
integer $E_n$. After deleting a finite prefix, all hypotheses of
Theorem~\ref{res:bounded} hold. Apply Corollary~\ref{res:cor}.
\end{proof}

The same calculation is valid with $qP_n$ replaced by any positive
clearance $H_n\le qP_n$: if $W_n=H_nx_n$ and
$Z_n=H_n-(a_n-1)W_n$, then
\begin{equation}\label{eq:general-clearance-dictionary}
 Z_n+\frac{H_n}{a_n}\gamma_n=o(1).
\end{equation}
This will supply the LCM-prefactor consequence below.

\section{A scalar finite-mass criterion}\label{sec:mass}
\begin{theorem}[finite negative relative mass]\label{res:massscalar}\label{res:mass}
Let $C_n$ be positive integers and $E_n$ integers satisfying
$C_{n+1}=C_n-E_n$. If
\[
 \sum_n\frac{(-E_n)_+}{C_n}<\infty,
\]
then $E_n=0$ eventually. No denominator dynamics or normalised vanishing
is required.
\end{theorem}
\begin{proof}
Set $\delta_n=(-E_n)_+/C_n$. Then
\[
 C_N\le C_0\prod_{n<N}(1+\delta_n)
 \le C_0\exp\!\left(\sum_n\delta_n\right).
\]
Choose an integer upper bound $K$ for $C_n$. Each strict rise contributes
at least $1/K$ to $\sum\delta_n$, so there are only finitely many rises.
The remaining positive integer sequence is nonincreasing and stabilises.
\end{proof}

On an exact reciprocal-tail orbit the conclusion gives the Sylvester
recurrence. The scalar criterion and bounded negative error are distinct
quantities to estimate; on the common exact-orbit class with normalised
vanishing they are each equivalent to the same endpoint.

\section{Weighted first crossings of LCM records}\label{sec:lcmrecords}

Only steps setting an LCM numerator record need contribute. The first-crossing
argument permits subtraction of a fixed baseline and any decreasing weight
with divergent integral. The theorem in this section has an ordinary proof.

Set
\[
 L_n=\operatorname{lcm}(q,a_0,\ldots,a_{n-1}),\quad
 M_n=D_n/L_n,\quad U_n=C_n/M_n,\quad V_n=E_n/M_n.
\]
The rational tail has denominator dividing $L_n$, so $U_n$ and $V_n$ are
integers. With $\rho_n=\gcd(L_n,a_n)$, the exact updates are
\[
 M_{n+1}=M_n\rho_n,\qquad
 \rho_nU_{n+1}=U_n-V_n,\qquad V_n=L_n-(a_n-1)U_n.
\]
Write $R_n=\max_{j\le n}U_j$ and
$\mathcal R=\{n:U_{n+1}>R_n\}$. Centring implies that every sufficiently
late strict rise has $\rho_n=1$: otherwise $U_{n+1}\le3U_n/4$.
Its actual jump is therefore $d_n=U_{n+1}-U_n=-V_n$.

\begin{theorem}[weighted record excess]\label{res:weightedrecord}
Assume the growth and rationality hypotheses of Problem~\ref{res:problem}.
Let $f:[1,\infty)\to[0,\infty)$ be finite and nonincreasing, with
$\int_1^\infty f(t)\,dt=\infty$. Then the sequence is eventually Sylvester
if and only if, for some integer $B\ge0$,
\[
 \sum_{n\in\mathcal R}(-V_n-B)_+f(U_n)<\infty.
\]
\end{theorem}

\begin{proof}
Suppose the sequence is not eventually Sylvester. Absorption and
integrality give $1/U_n\le |V_n|/U_n=|E_n|/C_n\to0$, so there are
infinitely many record steps. Every sufficiently late record is fresh.
Their digits are pairwise coprime: an earlier digit divides the later
$L_n$, while the digit at a fresh record is coprime to $L_n$.
For any fixed $B\ge1$, choose $B$ such digits
$m_0,\ldots,m_{B-1}>B$ and take $T$ after their indices.

Put $P=\prod_i m_i$ and choose $x$ by the Chinese remainder theorem with
$m_i\mid x+i$. Consider all translates $\tau=x+B+kP>R_T$, $k\in\Z$;
the first is at most $R_T+P$. A first crossing
$U_n\le R_n<\tau\le U_n+d_n$ is a record step. If $d_n\le B$, then
$U_n\in[\tau-B,\tau)$, so some $m_i$ divides $U_n$. It also divides
$L_n$, hence divides $d_n=(a_n-1)U_n-L_n$, contradicting
$0<d_n\le B<m_i$.

If the step first crosses $h\ge1$ such heights, their spacing gives
$(h-1)P<d_n$. With $r=d_n-B\ge1$ and $P\ge B+1$, we have
$d_n=B+r\le Pr$, whence $h\le r$. Monotonicity of $f$ now gives
\[
 \sum_{\substack{\tau\text{ first crossed}\\\text{at step }n}}f(\tau)
 \le(d_n-B)f(U_n).
\]
Each height above $R_T$ has one first crossing. Summing, and comparing each
interval of length $P$ with its left endpoint, yields the finite bound
\begin{equation}\label{eq:weightedcrossing}
 \sum_{\substack{T\le n<N\\n\in\mathcal R}}(-V_n-B)_+f(U_n)
 \ge\frac1P\int_{R_T+P}^{R_N} f(t)\,dt.
\end{equation}
Since $R_n\to\infty$, the right side diverges for every $B\ge1$;
$B=0$ follows by domination. Conversely a Sylvester tail telescopes to
$x_n=1/(a_n-1)$, so $V_n=0$ eventually.
\end{proof}

For example, $f(t)=1/[t\log(et)]$ gives the sufficient condition
\[
 \sum_{n\in\mathcal R}\frac{(-V_n-B)_+}{U_n\log(eU_n)}<\infty.
\]
Its finite lower bound in~\eqref{eq:weightedcrossing} is
$P^{-1}\log\bigl(\log(eR_N)/\log(e(R_T+P))\bigr)$.
Further fixed iterated logarithmic factors are allowed whenever the integral
still diverges. These are specialisations of one crossing theorem.

The criterion also has an exact expression in the original growth defect.
Put $\gamma_n=a_n^2/a_{n+1}-1$ and $\theta_n=E_n/C_n$. The defect identity
gives
\[
 \gamma_n+\theta_n=
 \frac{(1-\theta_n)(a_n-1+\theta_{n+1})}{a_{n+1}},
 \qquad 0<\gamma_n+\theta_n<3/a_n
\]
eventually. Thus the two nonnegative summands
$U_nf(U_n)(\gamma_n-B/U_n)_+$ and $(-V_n-B)_+f(U_n)$ differ by at most
$3U_nf(U_n)/a_n$. This is summable because
$U_n\le C_n=\exp(o(n))$, $f(U_n)\le f(1)$, and
$a_n\ge\exp(c2^n)$ eventually for some $c>0$.
Consequently Theorem~\ref{res:weightedrecord} is equivalent to finiteness of
\begin{equation}\label{eq:weightedgrowth}
 \sum_{n\in\mathcal R}U_nf(U_n)
 \left(\frac{a_n^2}{a_{n+1}}-1-\frac B{U_n}\right)_+
\end{equation}
for some $B$. The original hypotheses do not currently supply this
finiteness. In particular, termwise convergence to zero is insufficient.
Also, $d_n$ is the actual jump, including any recovery from a drawdown; it
must not be replaced by $R_{n+1}-R_n$ in the crossing proof.

\begin{corollary}[LCM-weighted bounded defect]\label{res:lcmbounded}
Assume the hypotheses of Problem~\ref{res:problem}. Write
$A_n=\operatorname{lcm}(a_1,\ldots,a_{n-1})$ with $A_1=1$ and
\[
 Z_n=\frac{A_n}{a_n}\Bigl(\frac{a_n^2}{a_{n+1}}-1\Bigr).
\]
If $\limsup Z_n<\infty$, then the sequence is eventually Sylvester.
\end{corollary}

\begin{proof}
Put $t_n=L_n/A_n=q/\gcd(q,A_n)$, so $1\le t_n\le q$.
Equation~\eqref{eq:general-clearance-dictionary}, with $H_n=L_n$, gives
$V_n+t_nZ_n=o(1)$. A finite upper bound on $Z_n$ therefore bounds the
negative part of $V_n$. The record series vanishes after a finite prefix
for a sufficiently large baseline, so Theorem~\ref{res:weightedrecord}
applies.
\end{proof}

Erd\H{o}s--Straus Theorem~3 uses a nonpositive limsup in the corresponding
next-index LCM expression \cite{erdosstraus1964}. Here any finite upper
bound suffices under the quadratic-limit assumption. Since $A_n\mid P_n$,
this also implies Corollary~\ref{res:originalbounded}: multiplication by
a factor in $(0,1]$ preserves an upper bound, including at negative values.

The divisor constraint also survives integral numerator coefficients:
$d=(a_n-1)U_n-b_nL_n$ is divisible by every divisor of both $U_n,L_n$.
The already-proved coefficient-uniform variant gives bounded LCM height
from a bounded negative part; normalised vanishing is used afterwards for
stationarity. Its proof and exact bounded-height counterexamples are in the
\href{\repobase/\PX/Erdos243/CoefficientUniformBoundedHeight.md}
{coefficient proof supplement}.

\section{Further consequences}\label{sec:secondaryrate}
\begin{theorem}[cubic-rate irrationality]\label{res:cubicrate}
A strictly increasing sequence of positive integers with
\[
 a_n^2/a_{n+1}=1+\frac3n+o(n^{-3})
\]
has irrational reciprocal sum.
\end{theorem}
\begin{proof}[Proof by the polynomial exclusion]
Under rationality, the canonical estimate gives
$C_{n+1}/C_n=1+3/n+o(n^{-3})$. Integer finite differences force
$C_n=An(n+1)(n+2)+B$ eventually. The fixed-cubic exclusion rules this out.
The complete ordinary polynomial and number-field argument is retained in
\href{\repobase/\PX/Erdos243/WeightedRecordOctupleCriteria.md\#8-rising-factorial-cubic-profiles-and-the-cubic-rate-irrationality-r08}
{the proof supplement, Section~8}; finite congruence checks alone are
insufficient for that implication.
\end{proof}

The same finite-difference extraction excludes every nonintegral
$\lambda>1$ under the rate
$a_n^2/a_{n+1}=1+\lambda/n+o(n^{-\lambda})$: eventual polynomial growth
would force its degree to equal $\lambda$. These rates lie outside
Koizumi's $1+o(1/n)$ hypothesis \cite[Remark~3]{koizumi2025}.

There is also an admitted quantitative extension. Under the hypotheses of
Theorem~\ref{res:bounded} other than (5), put
$\operatorname{LL}(x)=\log_2\log_2\max(4,x)$. Each fixed
$\delta\in(0,1)$ gives
\[
 (-E_n)_+\le(1-\delta)\operatorname{LL}(C_n)
 \quad\text{eventually}\quad\Longrightarrow\quad E_n=0\text{ eventually}.
\]
The LCM version needs the corresponding bound only at late record steps.
The proofs use a CRT block in $[P,2P)$ and the canonical multiplier scale;
they are retained in
\href{\repobase/\PX/Erdos243/SlowNegativePartRigidity.md}{the slow-negative proof}
and \href{\repobase/\PX/Erdos243/LcmRecordExcess.md\#5-record-only-subcritical-log-log-bound}
{the record-only extension}. These are ordinary results with stated rate
hypotheses; neither establishes the general record budget.

\section{The remaining arithmetic estimate}\label{sec:open}
\begin{proposition}[necessary profile]\label{res:frontier}
The canonical state of a counterexample has $E_n\ne0$ eventually,
$|E_n|/C_n\to0$, unbounded negative magnitudes along negative indices,
and
\[
 \sum_n\frac{(-E_n)_+}{C_n}=\infty.
\]
\end{proposition}
\begin{proof}
Absorption excludes late zeros, descent excludes an eventually
nonnegative error, and Theorems~\ref{res:bounded} and~\ref{res:massscalar}
exclude the two finiteness conditions.
\end{proof}

For $B\ge0$ put
\[
 F_B(X)=\sum_{\substack{n\in\mathcal R\\U_n\le X}}(-V_n-B)_+.
\]
The remaining equivalent estimate is
\begin{equation}\label{eq:remaining-record-budget}
 \exists B\in\N:\qquad\liminf_{X\to\infty}\frac{F_B(X)}X=0.
\end{equation}
On a nonterminal orbit first crossings instead give
$F_B(X)\ge X/P_B-O_B(1)$, with an orbit-dependent CRT modulus $P_B$.
The missing step is to derive~\eqref{eq:remaining-record-budget} from
the canonical dynamics. The equivalence with the existence of an
admissible summable weight is proved in the
\href{\repobase/\PX/Erdos243/LcmDefectCriterionReduction.md}
{record-budget supplement}; it does not establish the estimate.

The obstruction is the exact unit-numerator feedback at the record sources.
The scalar sequence $C_n=n^2+1$, $E_n=-(2n+1)$ has normalised vanishing,
subexponential height and finite square mass, but is not an exact
reciprocal-tail orbit: its numerator word $0,0,2\pmod5$ violates the
persistent-zero transport of the exact equations. Likewise, small rises
and avoidance of an arbitrary sparse coprime family do not control the
canonical timing and size of the available divisors. The supporting record
keeps those falsifying examples with the exact hypotheses each preserves.

\paragraph{Formal scope and supporting record.}
The bounded-negative state theorem is checked in Lean 4 against Mathlib
at the stated source checkpoint:
\lword{Erdos243/ReciprocalTailRigidity.lean}{2352}{eventuallyBoundedNegativePart_eventually_zero}{eventual bounded-negative rigidity}.
The canonical analytic transfer and the global weighted, polynomial and
slow-negative arguments are ordinary proofs; the scalar finite-mass proof
is printed in full above. The auxiliary source guide, constant and periodic
exclusions, primitive feedback, finite certificates and unsuccessful global
extensions belong to the accompanying reasoning record. Their stronger
global producer hypotheses remain explicit.

\paragraph{Declarations.}
The authorship and AI-use statement on the first page applies to this
manuscript. Code and mathematical source are available at the repository
checkpoint \texttt{\commitshort}. A source declaration establishes only
its own statement; the ordinary transfer and the unresolved estimate above
are separately identified.

\begin{thebibliography}{9}
\small
\raggedright
\setlength{\itemsep}{0pt}
\bibitem{erdosstraus1964} P.~Erd\H{o}s and E.~G. Straus,
  \href{https://users.renyi.hu/~p_erdos/1964-19.pdf}{\emph{On the irrationality
  of certain Ahmes series}},
  J. Indian Math. Soc. (N.S.) \textbf{27} (1964), 129--133. MR~175848.
\bibitem{duverney2001} D.~Duverney,
  \href{https://www.ms.u-tokyo.ac.jp/journal/pdf/jms080206.pdf}
  {\emph{Irrationality of fast converging series of rational numbers}},
  J. Math. Sci. Univ. Tokyo \textbf{8} (2001), 275--316. MR~1837165.
\bibitem{badea1993} C.~Badea,
  \href{https://matwbn.icm.edu.pl/ksiazki/aa/aa63/aa6342.pdf}{\emph{A theorem
  on irrationality of infinite series and applications}}, Acta Arith.
  \textbf{63} (1993), no.~4, 313--323.
\bibitem{erdosgraham1980} P.~Erd\H{o}s and R.~L. Graham,
  \href{https://mathweb.ucsd.edu/~ronspubs/80_11_number_theory.pdf}{\emph{Old
  and New Problems and Results in Combinatorial Number Theory}},
  Monogr. Enseign. Math. 28, Geneva, 1980, p.~64.
\bibitem{erdos1988} P.~Erd\H{o}s, \emph{On the irrationality of certain series:
  problems and results}, in A.~Baker (ed.), \emph{New Advances in Transcendence
  Theory}, Cambridge UP, 1988, pp.~102--109,
  doi:\href{https://doi.org/10.1017/CBO9780511897184.009}{10.1017/CBO9780511897184.009}.
\bibitem{kovactao2024} V.~Kova\v{c} and T.~Tao,
  \href{https://arxiv.org/abs/2406.17593}{\emph{On several irrationality
  problems for Ahmes series}}, Acta Math. Hungar. \textbf{175} (2025),
  572--608; preprint arXiv:2406.17593v4.
\bibitem{bado2026} I.~O. Bado,
  \emph{Prime-Support Rigidity and Primitive Pseudo-Greedy Dynamics:
  Partial Progress on Erd\H{o}s Problem~\#243}, author-posted preprint,
  September 2026, doi:\href{https://doi.org/10.13140/RG.2.2.36612.08325}
  {10.13140/RG.2.2.36612.08325}.  Not independently validated here.
\bibitem{koizumi2025} J.~Koizumi,
  \href{https://math.colgate.edu/~integers/aa28/aa28.pdf}
  {\emph{Irrationality of the reciprocal sum of doubly exponential sequences}},
  Integers \textbf{26} (2026), Paper No.~A28, 17 pp.,
  doi:\href{https://doi.org/10.5281/zenodo.18714404}{10.5281/zenodo.18714404};
  preprint arXiv:2504.05933v1.
\bibitem{erdosproblems} T.~F. Bloom,
  \href{https://www.erdosproblems.com/243}{\emph{Erd\H{o}s Problem \#243}},
  \texttt{erdosproblems.com/243}.
\end{thebibliography}

\end{document}
